\documentclass[aspectratio=169,10pt]{beamer}

\usetheme{Madrid}
\usecolortheme{default}
\usepackage[T1]{fontenc}
\usepackage{lmodern}
\usepackage{amsmath}
\usepackage{booktabs}
\usepackage{array}
\usepackage{listings}
\usepackage{xcolor}

\definecolor{codeblue}{RGB}{30,75,135}
\definecolor{codegreen}{RGB}{0,105,70}
\definecolor{codegray}{RGB}{95,95,95}
\lstset{
  basicstyle=\ttfamily\scriptsize,
  keywordstyle=\color{codeblue}\bfseries,
  commentstyle=\color{codegreen},
  stringstyle=\color{codegray},
  showstringspaces=false,
  breaklines=true,
  frame=single,
  columns=fullflexible
}

\newcommand{\shape}[1]{\texttt{#1}}
\newcommand{\file}[1]{\texttt{#1}}
\newcommand{\code}[1]{\texttt{#1}}

\title{Hybrid Convolutional Transformer\\for Face-Image Forensic Hashing}
\subtitle{Architecture and top-level code walkthrough}
\author{Convolutional\_Transformer repository}
\date{\today}

\begin{document}

\begin{frame}
  \titlepage
\end{frame}

\begin{frame}{Scope of this walkthrough}
  \begin{itemize}
    \item Covers every top-level source file in the repository root.
    \item Python: \file{CNN.py}, \file{tokenizer.py}, \file{transformer.py},
          \file{augmentations.py}, \file{loss.py}, \file{train.py}, \file{eval.py}.
    \item HPC launch scripts: \file{submit\_CvT\_train.sh}, \file{submit\_4gpu.sh}.
    \item Sub-folders are intentionally excluded.
    \item Focus: exact tensor shapes, kernel sizes, layer counts, optimization
          behavior, distributed training, evaluation, and implementation caveats.
  \end{itemize}
\end{frame}

\begin{frame}{Problem formulation}
  \begin{block}{Goal}
    Learn a semantic forensic hash that is stable under benign global image
    degradation but changes when a localized manipulation is introduced.
  \end{block}
  \begin{columns}[T]
    \column{0.32\textwidth}
      \textbf{Anchor}\\
      Pristine face image
    \column{0.32\textwidth}
      \textbf{Positive view}\\
      Global benign blur, noise, and brightness shift
    \column{0.32\textwidth}
      \textbf{Hard negative}\\
      Localized heavily blurred rectangular region
  \end{columns}
  \vspace{0.7em}
  Every view is encoded into a normalized 2,048-dimensional vector. Training
  pulls the anchor toward the benign view and pushes it away from the forged view
  and unrelated images.
\end{frame}

\begin{frame}{End-to-end architecture}
  \centering
  \scriptsize
  \setlength{\tabcolsep}{1pt}
  \begin{tabular}{>{\centering\arraybackslash}p{0.13\textwidth}
                  >{\centering\arraybackslash}p{0.04\textwidth}
                  >{\centering\arraybackslash}p{0.16\textwidth}
                  >{\centering\arraybackslash}p{0.04\textwidth}
                  >{\centering\arraybackslash}p{0.16\textwidth}
                  >{\centering\arraybackslash}p{0.04\textwidth}
                  >{\centering\arraybackslash}p{0.15\textwidth}
                  >{\centering\arraybackslash}p{0.04\textwidth}
                  >{\centering\arraybackslash}p{0.14\textwidth}}
    RGB image & $\rightarrow$ & Truncated ResNet-50 & $\rightarrow$ &
    Spatial tokenizer & $\rightarrow$ & 6 Transformer blocks & $\rightarrow$ &
    Semantic projector\\[0.5em]
    \shape{[B,3,512,512]} && \shape{[B,1024,32,32]} &&
    \shape{[B,1025,768]} && \shape{[B,1025,768]} &&
    \shape{[B,2048]}
  \end{tabular}
  \vspace{1em}
  \begin{itemize}
    \item The CNN preserves local forensic texture information on a $32\times32$ grid.
    \item The tokenizer turns each grid cell into one token and prepends a learnable
          classification token.
    \item Self-attention allows every spatial location to interact with every other location.
    \item The final classification token becomes the compact semantic hash.
  \end{itemize}
\end{frame}

\begin{frame}{Parameter budget}
  \centering
  \begin{tabular}{lr}
    \toprule
    Module & Trainable parameters\\
    \midrule
    Truncated ResNet-50 backbone & 8,543,296\\
    Spatial tokenizer & 1,575,168\\
    Six transformer blocks & 42,527,232\\
    Semantic projector & 1,579,008\\
    \midrule
    \textbf{Total} & \textbf{54,224,704}\\
    \bottomrule
  \end{tabular}
  \vspace{0.8em}
  \begin{itemize}
    \item One transformer block has 7,087,872 trainable parameters.
    \item The six-block transformer head contributes about 78.4\% of the total.
    \item Augmentation kernels are registered buffers, not trainable parameters.
  \end{itemize}
\end{frame}

\section{CNN Backbone}

\begin{frame}{\file{CNN.py}: truncated ResNet-50}
  \begin{itemize}
    \item Implements ResNet bottleneck blocks from scratch.
    \item Uses the standard ResNet-50 stage depths $[3,4,6,3]$, but omits the
          last stage: only $[3,4,6]$ are constructed.
    \item Also omits adaptive average pooling and the classification layer.
    \item Input is resized to $512\times512$ before it reaches the backbone.
    \item Output is a spatial map, not a class score:
          \shape{[B,1024,32,32]}.
  \end{itemize}
  \begin{block}{Why truncate?}
    Keeping the $32\times32$ map supplies 1,024 spatial tokens. A complete
    ResNet-50 would downsample further and discard useful local anomaly detail.
  \end{block}
\end{frame}

\begin{frame}{CNN stem: exact downsampling}
  \centering
  \begin{tabular}{lllll}
    \toprule
    Operation & Kernel & Stride & Padding & Output shape\\
    \midrule
    Input & -- & -- & -- & \shape{[B,3,512,512]}\\
    Conv2d & $7\times7$ & 2 & 3 & \shape{[B,64,256,256]}\\
    BatchNorm2d + ReLU & -- & -- & -- & \shape{[B,64,256,256]}\\
    MaxPool2d & $3\times3$ & 2 & 1 & \shape{[B,64,128,128]}\\
    \bottomrule
  \end{tabular}
  \vspace{0.8em}
  For convolution or pooling, the spatial output is
  \[
    \left\lfloor\frac{H + 2P - K}{S}\right\rfloor + 1.
  \]
  All convolutions use \code{bias=False} when immediately followed by batch normalization.
\end{frame}

\begin{frame}{Residual bottleneck block}
  Each \code{Bottleneck} applies:
  \[
    x \rightarrow
    \underbrace{\mathrm{Conv}_{1\times1}}_{\text{reduce}}
    \rightarrow \mathrm{BN}\rightarrow\mathrm{ReLU}
    \rightarrow
    \underbrace{\mathrm{Conv}_{3\times3}}_{\text{spatial}}
    \rightarrow \mathrm{BN}\rightarrow\mathrm{ReLU}
    \rightarrow
    \underbrace{\mathrm{Conv}_{1\times1}}_{\text{expand by 4}}
    \rightarrow \mathrm{BN}.
  \]
  Then:
  \[
    y = \mathrm{ReLU}(F(x) + D(x)).
  \]
  \begin{itemize}
    \item Expansion factor: 4.
    \item The $3\times3$ convolution uses padding 1 and carries the stage stride.
    \item $D(x)$ is either the identity or a $1\times1$ projection plus BatchNorm.
    \item ReLU is configured with \code{inplace=False}.
  \end{itemize}
\end{frame}

\begin{frame}{CNN stages: dimensions and block counts}
  \centering
  \scriptsize
  \begin{tabular}{lllllll}
    \toprule
    Stage & Blocks & Internal channels & Expanded output & First-block stride & Grid & Parameters\\
    \midrule
    Stem & -- & -- & 64 & 2 then pool 2 & $128\times128$ & 9,536\\
    Layer 1 & 3 & 64 & 256 & 1 & $128\times128$ & 215,808\\
    Layer 2 & 4 & 128 & 512 & 2 & $64\times64$ & 1,219,584\\
    Layer 3 & 6 & 256 & 1,024 & 2 & $32\times32$ & 7,098,368\\
    \midrule
    Output & -- & -- & 1,024 & -- & $32\times32$ & 8,543,296 total\\
    \bottomrule
  \end{tabular}
  \vspace{0.7em}
  A projection shortcut is created whenever stride changes or channel counts do
  not match the expanded output. CNN convolutions use Kaiming-normal
  initialization; BatchNorm scales start at 1 and biases at 0.
\end{frame}

\section{Tokenizer}

\begin{frame}{\file{tokenizer.py}: spatial-to-sequence bridge}
  The tokenizer preserves all $32\times32=1,024$ CNN grid positions:
  \[
    \shape{[B,1024,32,32]}
    \rightarrow
    \shape{[B,1024,1024]}
    \rightarrow
    \shape{[B,1024,768]}.
  \]
  \begin{enumerate}
    \item \code{view} flattens the spatial axes.
    \item \code{permute(0,2,1)} changes channel-first CNN layout into token-first layout.
    \item A learned linear projection maps each 1,024-channel CNN vector to 768 dimensions.
    \item A learnable \shape{[1,1,768]} classification token is prepended.
    \item A learnable \shape{[1,1025,768]} absolute positional embedding is added.
  \end{enumerate}
\end{frame}

\begin{frame}{Tokenizer parameters and initialization}
  \centering
  \begin{tabular}{lll}
    \toprule
    Object & Shape & Parameters\\
    \midrule
    Patch projection weight & \shape{[768,1024]} & 786,432\\
    Patch projection bias & \shape{[768]} & 768\\
    Classification token & \shape{[1,1,768]} & 768\\
    Positional embedding & \shape{[1,1025,768]} & 787,200\\
    \midrule
    Total & -- & 1,575,168\\
    \bottomrule
  \end{tabular}
  \vspace{0.8em}
  \begin{itemize}
    \item Classification token and positional embedding: truncated normal, standard deviation 0.02.
    \item Projection weight: Xavier-uniform.
    \item Projection bias: zeros.
    \item There is no patchifying convolution: tokens correspond directly to CNN grid cells.
  \end{itemize}
\end{frame}

\section{Transformer}

\begin{frame}{\file{transformer.py}: six pre-norm blocks}
  The model creates six identical \code{CustomTransformerBlock} instances:
  \[
    Z_{\ell}' = Z_{\ell} + \mathrm{MHSA}(\mathrm{LN}(Z_{\ell})),
  \]
  \[
    Z_{\ell+1} = Z_{\ell}' + \mathrm{MLP}(\mathrm{LN}(Z_{\ell}')).
  \]
  \begin{itemize}
    \item Hidden dimension: 768.
    \item Sequence length: 1,025.
    \item Attention heads: 12.
    \item Per-head dimension: $768/12=64$.
    \item MLP ratio: 4, giving a 3,072-dimensional hidden layer.
    \item Both sublayers have residual additions.
  \end{itemize}
\end{frame}

\begin{frame}{Multi-head self-attention in detail}
  \begin{enumerate}
    \item One fused linear layer projects \shape{[B,1025,768]} to
          \shape{[B,1025,2304]}.
    \item \code{chunk(3, dim=-1)} produces query, key, and value tensors.
    \item Each tensor is reshaped to \shape{[B,12,1025,64]}.
    \item PyTorch \code{scaled\_dot\_product\_attention} computes full,
          bidirectional self-attention with no mask and no dropout.
    \item Heads are concatenated back to \shape{[B,1025,768]}.
    \item A final linear layer mixes the concatenated heads.
  \end{enumerate}
  The mathematical operation per head is:
  \[
    \mathrm{Attention}(Q,K,V)=
    \mathrm{softmax}\left(\frac{QK^\top}{\sqrt{64}}\right)V.
  \]
\end{frame}

\begin{frame}{Attention implementation details}
  \begin{itemize}
    \item \code{attn\_mask=None}: all tokens can attend to all tokens.
    \item \code{is\_causal=False}: this is an encoder, not an autoregressive decoder.
    \item \code{dropout\_p=0.0}: attention is deterministic even in training.
    \item \code{scale=None}: PyTorch uses the default $1/\sqrt{d_k}$ factor.
    \item On a compatible CUDA stack and GPU, PyTorch may dispatch the SDPA call
          to an optimized Flash Attention backend.
    \item The code requests SDPA; the actual selected backend depends on runtime
          hardware, data type, and PyTorch configuration.
  \end{itemize}
\end{frame}

\begin{frame}{MLP and optional checkpointing}
  The second transformer sublayer is:
  \[
    \shape{768}
    \xrightarrow{\mathrm{Linear}}
    \shape{3072}
    \xrightarrow{\mathrm{GELU}}
    \shape{3072}
    \xrightarrow{\mathrm{Linear}}
    \shape{768}.
  \]
  \begin{itemize}
    \item Each block exposes \code{use\_checkpoint=False} by default.
    \item Training can set it to true for every block via
          \code{--enable-checkpointing}.
    \item During training, \code{checkpoint(..., use\_reentrant=False)}
          recomputes block activations during backward to reduce activation memory.
    \item The trade-off is extra computation during backpropagation.
  \end{itemize}
\end{frame}

\begin{frame}{One transformer block: parameter accounting}
  \centering
  \scriptsize
  \begin{tabular}{lr}
    \toprule
    Component & Parameters\\
    \midrule
    Two LayerNorm layers & 3,072\\
    Fused QKV projection: $768\rightarrow2304$ & 1,771,776\\
    Attention output projection: $768\rightarrow768$ & 590,592\\
    MLP first linear: $768\rightarrow3072$ & 2,362,368\\
    MLP second linear: $3072\rightarrow768$ & 2,360,064\\
    \midrule
    One block & 7,087,872\\
    Six blocks & 42,527,232\\
    \bottomrule
  \end{tabular}
\end{frame}

\begin{frame}{Semantic projector}
  After the sixth transformer block, only token index 0 is retained:
  \[
    \shape{[B,1025,768]}
    \xrightarrow{\text{select CLS}}
    \shape{[B,768]}
    \xrightarrow{\mathrm{Linear}}
    \shape{[B,2048]}
    \xrightarrow{\mathrm{GELU + LayerNorm}}
    \shape{[B,2048]}.
  \]
  \begin{itemize}
    \item Linear layer: $768\rightarrow2,048$.
    \item Projector parameters: 1,579,008.
    \item The loss later L2-normalizes hashes onto the unit hypersphere.
    \item The model itself returns the unnormalized 2,048-D vector.
  \end{itemize}
\end{frame}

\section{Augmentations}

\begin{frame}{\file{augmentations.py}: fixed depthwise filter bank}
  Raw filters are non-trainable registered buffers. Each RGB channel is filtered
  independently with \code{groups=3}.
  \vspace{0.5em}
  \centering
  \begin{tabular}{llll}
    \toprule
    Filter & Kernel & Default parameters & Purpose\\
    \midrule
    \code{RawGaussianFilter} & $K\times K$ & $K=7,\sigma=1.5$ & Gaussian blur\\
    \code{RawBoxFilter} & $K\times K$ & $K=5$ & Uniform averaging\\
    \code{RawMotionFilter} & $K\times K$ & $K=7$ & Diagonal camera shake\\
    \bottomrule
  \end{tabular}
  \vspace{0.6em}
  For odd $K$, padding is $K//2$, preserving image height and width. Each
  kernel is reshaped to \shape{[1,1,K,K]} and repeated to
  \shape{[3,1,K,K]}.
\end{frame}

\begin{frame}{Gaussian, box, and motion kernels}
  \begin{block}{Gaussian kernel}
    Coordinates run from $-P$ to $P$, where $P=K//2$. The 1-D Gaussian is
    $g(x)=\exp(-x^2/(2\sigma^2))$. The 2-D kernel is the normalized outer
    product $gg^\top$.
  \end{block}
  \begin{block}{Box kernel}
    Every one of the $K^2$ entries is $1/K^2$, so brightness is preserved.
  \end{block}
  \begin{block}{Motion kernel}
    A normalized $K\times K$ identity matrix places $1/K$ along the diagonal,
    producing an anisotropic diagonal smear.
  \end{block}
\end{frame}

\begin{frame}{Global benign augmentation}
  For an entire batch, \code{GlobalBenignAugmentation} randomly selects one blur:
  \begin{itemize}
    \item Gaussian: $5\times5$, $\sigma=1.0$.
    \item Gaussian: $7\times7$, $\sigma=2.0$.
    \item Box: $5\times5$.
    \item Motion: $7\times7$.
  \end{itemize}
  It then adds pixelwise Gaussian noise with standard deviation sampled uniformly
  from $[0.01,0.05]$, multiplies by a brightness scalar sampled uniformly from
  $[0.9,1.1]$, and clamps pixels to $[0,1]$.
  \begin{block}{Implementation detail}
    The bank is an \code{nn.ModuleList}; therefore \code{.to(device)} moves every
    registered filter buffer to the GPU.
  \end{block}
\end{frame}

\begin{frame}{Localized malicious forgery}
  Training constructs \code{LocalMaliciousForgery(tamper\_size=(96,192))}.
  \begin{enumerate}
    \item Pick one heavy Gaussian blur for the batch: $31\times31$ kernel with
          $\sigma\in\{5,6,7,8\}$.
    \item For each image, independently sample rectangle height and width from
          $[96,192]$ pixels.
    \item Sample a rectangle origin and set the binary mask region to 1.
    \item Blur the full image, then use the mask to blend blurred pixels only
          into the selected region.
  \end{enumerate}
  The class default is \code{(64,128)}, but \file{train.py} overrides it.
\end{frame}

\begin{frame}{Soft mask for forgery blending}
  A binary rectangle would create an obvious hard edge. The code feathers that
  edge by convolving the one-channel mask with a Gaussian:
  \[
    M_{\mathrm{soft}} = M_{\mathrm{hard}} * G_{31\times31,\sigma=7}.
  \]
  The final image is:
  \[
    x_{\mathrm{tampered}} =
    M_{\mathrm{soft}}\odot x_{\mathrm{blurred}} +
    (1-M_{\mathrm{soft}})\odot x.
  \]
  \begin{itemize}
    \item Soft-mask kernel: $31\times31$, $\sigma=7$, padding 15.
    \item One randomized box is made per image.
    \item Heavy-blur choice is shared by the batch.
    \item CUDA scalar indices are converted to Python integers with \code{.item()}.
  \end{itemize}
\end{frame}

\section{Loss}

\begin{frame}{\file{loss.py}: asymmetric contrastive objective}
  Let $h_o$, $h_g$, and $h_t$ denote pristine, globally augmented, and locally
  tampered hashes. First normalize:
  \[
    z = \frac{h}{\lVert h\rVert_2}.
  \]
  With temperature $\tau=0.05$ and hard-negative weight $\alpha=20$:
  \[
    s_i^+ = \frac{z_{o,i}^\top z_{g,i}}{\tau},
    \qquad
    s_i^{\mathrm{hard}} = \frac{z_{o,i}^\top z_{t,i}}{\tau}.
  \]
  Anchor direction is asymmetric: pristine images are anchors, benign images
  are positives, and tampered images are weighted hard negatives.
\end{frame}

\begin{frame}{Implemented denominator}
  For each anchor, compare against all gathered benign views and remove its
  matched positive from that set. The implemented loss is:
  \[
    \mathcal{L}_i =
    -\left[
      s_i^+ -
      \log\left(
        \sum_{j\ne i}\exp(s_{ij}) +
        \alpha\exp(s_i^{\mathrm{hard}})
      \right)
    \right].
  \]
  \begin{itemize}
    \item $s_{ij}=z_{o,i}^\top z_{g,j}/\tau$.
    \item Positive entries are replaced with $-\infty$ before exponentiation.
    \item The batch loss is the mean of per-sample losses.
    \item Similarity computation is cosine similarity because vectors are normalized.
  \end{itemize}
\end{frame}

\begin{frame}{Distributed negative gathering}
  In multi-GPU mode:
  \begin{itemize}
    \item Each rank gathers benign-view embeddings from all ranks.
    \item Local batch size $B$ becomes global context size
          $N_{\mathrm{global}}=B\times\mathrm{WORLD\_SIZE}$.
    \item \code{dist.all\_gather} does not preserve autograd for gathered tensors.
    \item The code replaces the gathered local shard with its original tensor,
          preserving gradients for that local shard.
    \item The local positive columns begin at \code{rank * B}; these positions
          are masked out of the denominator.
  \end{itemize}
  This expands the negative pool without sending gradients through remote shards.
\end{frame}

\begin{frame}{Loss formulation: presentation caveat}
  \begin{alertblock}{Verify the intended objective}
    The implementation excludes the positive term from the denominator. Standard
    InfoNCE / NT-Xent commonly includes the positive alongside negatives in the
    denominator. Excluding it can produce negative loss values when the positive
    similarity dominates.
  \end{alertblock}
  \begin{itemize}
    \item The slides describe the repository exactly as implemented.
    \item The code comments explicitly state that exclusion is intentional.
    \item Before comparing results against standard NT-Xent baselines, confirm
          whether the desired objective is standard InfoNCE or this custom ratio.
  \end{itemize}
\end{frame}

\section{Training}

\begin{frame}{\file{train.py}: dataset and data loading}
  \code{PristineFacesDataset}:
  \begin{itemize}
    \item Recursively discovers \code{.jpg}, \code{.jpeg}, and \code{.png} files.
    \item Converts each image to RGB.
    \item Resizes every image to $512\times512$.
    \item Converts pixels to a float tensor in $[0,1]$.
    \item Leaves augmentation to the GPU-side training loop.
  \end{itemize}
  \code{DataLoader}:
  \begin{itemize}
    \item Uses pinned memory, \code{drop\_last=True}, and prefetch factor 4.
    \item Persistent workers are enabled when worker count is positive.
    \item Uses \code{DistributedSampler} under DDP; otherwise uses shuffle.
  \end{itemize}
\end{frame}

\begin{frame}{Training batch construction}
  For each pristine mini-batch $x$:
  \[
    x_{\mathrm{global}} = A_{\mathrm{benign}}(x), \qquad
    x_{\mathrm{tampered}} = A_{\mathrm{local}}(x).
  \]
  Augmentations run inside \code{torch.no\_grad()}, then:
  \[
    x_{\mathrm{all}} =
    \mathrm{cat}(x,x_{\mathrm{global}},x_{\mathrm{tampered}})
    \in \mathbb{R}^{3B\times3\times512\times512}.
  \]
  A single model call returns \shape{[3B,2048]}, then \code{chunk(3)} recovers
  anchor, positive, and hard-negative hashes. This avoids three separate forward calls.
\end{frame}

\begin{frame}{Optimization configuration}
  \centering
  \begin{tabular}{ll}
    \toprule
    Setting & Default\\
    \midrule
    Optimizer & AdamW\\
    Base learning rate & $10^{-4}$\\
    Distributed learning rate & base LR $\times$ world size\\
    Weight decay & $10^{-4}$\\
    Epochs & 100 from Python CLI\\
    Warmup & 10 epochs\\
    Gradient clipping & max norm 1.0\\
    Temperature $\tau$ & 0.05\\
    Hard-negative weight $\alpha$ & 20.0\\
    Checkpoint interval & 10 epochs\\
    \bottomrule
  \end{tabular}
\end{frame}

\begin{frame}{Learning-rate schedule}
  \code{build\_lr\_lambda} applies linear warmup followed by cosine decay:
  \[
    \lambda(e)=
    \begin{cases}
      \dfrac{e+1}{E_{\mathrm{warmup}}}, & e < E_{\mathrm{warmup}},\\[0.8em]
      \dfrac{1}{2}\left(1+\cos\left(\pi
      \dfrac{e-E_{\mathrm{warmup}}}{E_{\mathrm{total}}-E_{\mathrm{warmup}}}
      \right)\right), & \text{otherwise.}
    \end{cases}
  \]
  \begin{itemize}
    \item Scheduler step occurs after each epoch.
    \item Learning rate is scaled linearly by the number of processes.
    \item A denominator guard of \code{max(1, ...)} handles edge cases.
  \end{itemize}
\end{frame}

\begin{frame}{CUDA, AMP, and DDP setup}
  \begin{itemize}
    \item Training requires CUDA; there is no CPU training fallback.
    \item If \code{LOCAL\_RANK}, \code{RANK}, and \code{WORLD\_SIZE} exist,
          NCCL is initialized. Partial distributed environments are rejected.
    \item Without those variables, the code uses one GPU for easier debugging.
    \item TF32 matmul and cuDNN TF32 are enabled.
    \item AMP selects bfloat16 when supported, otherwise float16.
    \item GradScaler is enabled only for float16.
    \item DDP uses one device per process and \code{find\_unused\_parameters=False}.
  \end{itemize}
\end{frame}

\begin{frame}{Checkpointing and resume behavior}
  After every epoch on rank 0, \file{train.py} saves:
  \begin{itemize}
    \item zero-based epoch index,
    \item unwrapped model state,
    \item optimizer state,
    \item scheduler state,
    \item averaged epoch loss.
  \end{itemize}
  The rolling path defaults to \file{latest\_checkpoint.pth}. Every tenth epoch,
  a named file such as \file{hybrid\_encoder\_epoch\_010.pth} is also written.
  If the rolling checkpoint already exists, training automatically resumes at
  the following epoch.
\end{frame}

\section{Evaluation}

\begin{frame}{\file{eval.py}: evaluation workflow}
  \file{eval.py} is an executable evaluation script with fixed constants:
  \begin{itemize}
    \item checkpoint: \file{hybrid\_encoder\_epoch\_150.pth},
    \item evaluation root: a hard-coded cluster path,
    \item device: CUDA when available, otherwise CPU,
    \item random seed: 42.
  \end{itemize}
  It loads the model, computes normalized hashes for \file{real},
  \file{positive}, and \file{negative} directories, and stores results in a
  directory named from the checkpoint epoch.
\end{frame}

\begin{frame}{Evaluation matrices and threshold search}
  If $R$, $P$, and $N$ contain normalized embedding rows, the script computes:
  \[
    M_{RP}=RP^\top,\qquad
    M_{RN}=RN^\top,\qquad
    M_{PN}=PN^\top.
  \]
  \begin{itemize}
    \item Diagonal entries are treated as matched image pairs.
    \item Off-diagonal entries are treated as unmatched pairs.
    \item Threshold candidates are sampled from $[-1,1]$ in increments of 0.001.
    \item A score above threshold predicts a benign positive pair.
    \item The threshold maximizing accuracy across matched RP and RN diagonals is saved.
  \end{itemize}
\end{frame}

\begin{frame}{Evaluation artifacts}
  The output directory contains:
  \begin{itemize}
    \item three serialized embedding tensors,
    \item \file{similarity\_matrices.npz},
    \item three 300-DPI heatmaps,
    \item \file{statistics.txt}.
  \end{itemize}
  The statistics report includes diagonal and off-diagonal mean, standard
  deviation, minimum, maximum, five random examples, the diagonal--off-diagonal
  gap, hard/easy real-negative cases, and best threshold accuracy.
\end{frame}

\section{HPC Scripts}

\begin{frame}{\file{submit\_CvT\_train.sh}: eight-GPU job}
  \centering
  \begin{tabular}{ll}
    \toprule
    SLURM resource & Requested value\\
    \midrule
    Nodes / tasks & 1 / 1\\
    GPUs & 8\\
    CPUs & 64\\
    Memory & 256 GB\\
    Wall time & 48 hours\\
    Partition / QoS & \code{dgx} / \code{dgx}\\
    \bottomrule
  \end{tabular}
  \vspace{0.7em}
  The script supports micromamba, mamba, or conda activation, launches
  \code{torchrun --nproc\_per\_node=8}, enables checkpointing, defaults to
  48 images per GPU and 500 epochs, and forwards additional command-line arguments.
\end{frame}

\begin{frame}{\file{submit\_4gpu.sh}: four-GPU job}
  \centering
  \begin{tabular}{ll}
    \toprule
    SLURM resource & Requested value\\
    \midrule
    Nodes / tasks & 1 / 1\\
    GPUs & 4\\
    CPUs & 32\\
    Memory & 128 GB\\
    Wall time & 24 hours\\
    Partition / QoS & \code{dgx} / \code{dgx}\\
    \bottomrule
  \end{tabular}
  \vspace{0.7em}
  This variant requires mamba, launches four ranks, enables checkpointing, and
  uses the same default 48 images per GPU and 500 epochs. Both scripts print
  environment details, invoke \code{nvidia-smi}, capture the exit code, and
  preserve strict shell error handling outside the training command.
\end{frame}

\begin{frame}{HPC environment variables}
  Both submission scripts configure:
  \begin{itemize}
    \item \code{OMP\_NUM\_THREADS=8},
    \item \code{NCCL\_DEBUG=INFO},
    \item \code{NCCL\_IB\_DISABLE=0},
    \item \code{NCCL\_P2P\_DISABLE=0},
    \item \code{PYTHONUNBUFFERED=1}.
  \end{itemize}
  User-overridable values include conda environment, project directory, dataset
  path, batch size, worker count, rolling checkpoint path, and epoch count.
  \begin{block}{Effective batches}
    Script defaults yield 192 images per step on four GPUs and 384 images per
    step on eight GPUs, before the training loop triples each batch into views.
  \end{block}
\end{frame}

\section{Code Notes}

\begin{frame}{Implementation notes to verify}
  \begin{itemize}
    \item \file{train.py} defaults to dataset path \file{Faces/train}; both shell
          scripts default to \file{Faces}. Recursive discovery may make either
          valid, but the intended layout should be documented.
    \item The Python CLI defaults to 100 epochs while both shell scripts override
          this with 500 epochs.
    \item \file{eval.py} accepts only top-level \file{*.jpg} files in each
          evaluation category; it ignores JPEG variants, PNG files, and nested files.
    \item \file{eval.py} has hard-coded checkpoint and dataset paths rather than CLI arguments.
    \item The contrastive denominator differs from standard InfoNCE / NT-Xent;
          confirm the intended research formulation.
  \end{itemize}
\end{frame}

\begin{frame}{Evaluation assumptions to verify}
  \begin{itemize}
    \item Similarity diagonals assume category folders use matching sorted filenames
          and compatible lengths.
    \item \code{compute\_stats} constructs a square identity mask from row count;
          non-square similarity matrices would fail or be interpreted incorrectly.
    \item Five random matched pairs require at least five samples.
    \item Off-diagonal sampling requires at least two samples.
    \item Empty image directories cause \code{torch.stack} to fail.
    \item Evaluation uses \code{torch.load} without \code{weights\_only=True};
          only trusted checkpoints should be loaded.
  \end{itemize}
\end{frame}

\begin{frame}{Training signal summary}
  \centering
  \begin{tabular}{p{0.28\textwidth}p{0.28\textwidth}p{0.28\textwidth}}
    \toprule
    Pristine anchor & Benign positive & Local hard negative\\
    \midrule
    Original resized RGB image &
    Global blur + noise + brightness shift &
    Feathered local heavy blur\\
    Should remain close to positive &
    Models acceptable acquisition changes &
    Should move away from anchor\\
    \bottomrule
  \end{tabular}
  \vspace{1em}
  \begin{block}{Core idea}
    CNN features preserve fine spatial evidence; transformer attention integrates
    evidence globally; contrastive training makes the final hash robust to benign
    degradation and sensitive to localized manipulation.
  \end{block}
\end{frame}

\begin{frame}{Complete tensor-shape trace}
  \centering
  \scriptsize
  \begin{tabular}{lll}
    \toprule
    Stage & Operation & Shape\\
    \midrule
    Input & RGB batch & \shape{[B,3,512,512]}\\
    CNN stem & $7\times7$ conv, $3\times3$ max pool & \shape{[B,64,128,128]}\\
    CNN layer 1 & 3 bottlenecks & \shape{[B,256,128,128]}\\
    CNN layer 2 & 4 bottlenecks & \shape{[B,512,64,64]}\\
    CNN layer 3 & 6 bottlenecks & \shape{[B,1024,32,32]}\\
    Tokenizer & flatten grid & \shape{[B,1024,1024]}\\
    Tokenizer & project channels & \shape{[B,1024,768]}\\
    Tokenizer & prepend CLS, add positions & \shape{[B,1025,768]}\\
    Transformer & 6 blocks & \shape{[B,1025,768]}\\
    Projector & select CLS & \shape{[B,768]}\\
    Projector & linear, GELU, LayerNorm & \shape{[B,2048]}\\
    Loss & L2 normalize & \shape{[B,2048]}\\
    \bottomrule
  \end{tabular}
\end{frame}

\begin{frame}{Top-level file map}
  \centering
  \begin{tabular}{ll}
    \toprule
    File & Responsibility\\
    \midrule
    \file{CNN.py} & Truncated ResNet-50 spatial backbone\\
    \file{tokenizer.py} & CNN grid to transformer-token sequence\\
    \file{transformer.py} & Attention blocks and semantic projector\\
    \file{augmentations.py} & Benign and malicious GPU-side image transforms\\
    \file{loss.py} & Asymmetric contrastive loss and distributed gather\\
    \file{train.py} & Dataset, model assembly, optimization, DDP, checkpoints\\
    \file{eval.py} & Embeddings, similarity matrices, thresholding, plots\\
    \file{submit\_CvT\_train.sh} & Eight-GPU SLURM launcher\\
    \file{submit\_4gpu.sh} & Four-GPU SLURM launcher\\
    \bottomrule
  \end{tabular}
\end{frame}

\begin{frame}
  \centering
  \Huge Questions?
\end{frame}

\end{document}
